\documentclass[xcolor=table,dvipsnames,table]{beamer}
\mode<presentation>
\usetheme{boxes}
\setbeamertemplate{navigation symbols}{}
% http://www.latex-community.org/forum/viewtopic.php?f=4&t=6694
\setbeamertemplate{navigation symbols}{\raisebox{5pt}{\makebox[\paperwidth]{\hfill\makebox[10pt]{\scriptsize\insertframenumber\vspace{1ex}}}}}
%\setbeamertemplate{footline}[frame number]
\setbeamertemplate{blocks}[shadow=false]
%\setbeamercolor*{block title}{fg=structure,bg=RoyalBlue!10}
\setbeamercolor*{block title example}{fg=structure,bg=RoyalBlue!10}
%\setbeamercolor*{block title example}{fg=BrickRed,bg=Goldenrod!10}
\setbeamercolor*{block title alerted}{fg=white,bg=black}
\addtobeamertemplate{block begin}{\pgfsetfillopacity{0.8}}{\pgfsetfillopacity{1}}
%\rowcolors{0}{RoyalBlue!20}{RoyalBlue!5}

%\DeclareGraphicsRule{*}{mps}{*}{}

\usepackage{latexsym}
\usepackage{hyperref}
\usepackage{tikz}
\usetikzlibrary{calc,shapes,arrows,shadows,shapes.callouts,shapes.arrows,chains,positioning,trees}
\usepackage{solution}
\usepackage{calc}
\usepackage{pifont}
\usepackage{algorithmic}

\newcommand{\cmark}{\ding{51}}
\newcommand{\xmark}{\ding{55}}

\newcounter{mycallout}

\newcommand{\callouts}[3]{%
  \stepcounter{mycallout}
  \tikz[remember picture,baseline]{\node[anchor=base,inner sep=0,outer sep=0]%
    (\themycallout) {\colorbox{#1!20}{#3}};\pause\node[overlay,rectangle callout,%
    callout relative pointer={(0cm,0.5cm)},fill=#1!20] at ($(\themycallout.south)+(-0cm,-0.7cm)$){#2};}%
    }%

\raggedright

\newcount\lecturecount
\lecturecount=0
\AtBeginLecture{%
    \advance\lecturecount by 1
    \date{}
    \begin{frame}
    \begin{center}
    \titlepage
    \ifnum\lecturecount=1
    Part \the\lecturecount: \insertlecture
    \else
    Part \the\lecturecount: \insertlecture
    \fi
    \end{center}
    \end{frame}
}

\addtobeamertemplate{block begin}{\setlength\abovedisplayskip{0pt}}

%\newcommand{\example}[1]{{\color{BrickRed!50}{#1}}}
\newcommand{\maths}[1]{{\color{RoyalBlue!50}{#1}}}
\newcommand{\reference}[1]{{\color{RoyalBlue!30}\tiny [from #1]}}
\newcommand{\koehnref}{\reference{\href{http://www.statmt.org/book}{P.Koehn SMT book slides}}}


\begin{document}

\title{\color{RoyalBlue!20}Natural Language Processing}

\author{Anoop Sarkar \\ {\tt anoopsarkar.github.io/nlp-class}}
\institute{Simon Fraser University}
%\date{}
     
{
\setbeamertemplate{navigation symbols}{}
\addtocounter{framenumber}{-1}
\begin{frame}
\begin{center}
\vspace{8mm}
\includegraphics[scale=0.35]{figures/natlang-cky-logo}
\end{center}
\titlepage
\end{frame}
}



\lecture{Linear models for Tagging}{}

\section{Tagging tasks in NLP}
\frame{\tableofcontents[currentsection]}
\newcommand{\postag}[1]{{\color{red}/#1}}
\newcommand{\nertag}[1]{{\color{blue}/#1}}

\begin{frame}
\frametitle{Tagging Tasks}
\begin{block}{Tagged Sequences}
a b e e a f h j $\Rightarrow$ a\postag{Y} b\postag{Z} e\postag{Y} e\postag{Y} a\postag{Z} f\postag{X} h\postag{Z} j\postag{Y}
\end{block}
\pause
\begin{alertblock}{Example 1: Part-of-speech tagging}
Profits\postag{N} soared\postag{V} at\postag{P} Boeing\postag{N} Co.\postag{N} ,\postag{,} easily\postag{ADV} topping\postag{V} forecasts\postag{N} on\postag{P} Wall\postag{N} Street\postag{N} ,\postag{,} as\postag{P} their\postag{POSS} CEO\postag{N} Alan\postag{N} Mulally\postag{N} announced\postag{V} first\postag{ADJ} quarter\postag{N} results\postag{N} .\postag{.}
\end{alertblock}
\pause
\begin{alertblock}{Example 2: Named Entity Recognition}
Profits\postag{O} soared\postag{O} at\postag{O} Boeing\nertag{B-CO} Co.\nertag{I-CO} ,\postag{O} easily\postag{O} topping\postag{O} forecasts\postag{O} on\postag{O} Wall\nertag{B-LOC} Street\nertag{I-LOC} ,\postag{O} as\postag{O} their\postag{O} CEO\postag{O} Alan\nertag{B-PER} Mulally\nertag{I-PER} announced\postag{O} first\postag{O} quarter\postag{O} results\postag{O} .\postag{O}
\end{alertblock}
\end{frame}


\begin{frame}
\frametitle{Notation for Tagging Tasks}
\begin{itemize}[<+->]
\item Set of possible input words: ${\cal V}$
\item Set of possible tags: ${\cal T}$
\item Word sequence: $x_{[1:n]} = [x_1, \ldots, x_n]$
\item Tag sequence: $t_{[1:n]} = [t_1, \ldots, t_n]$
\item Training data is $N$ tagged sentences, the $i^{th}$ sentence has length $n_i$:
\[ (x_{[1:n]}^{(i)}, t_{[1:n]}^{(i)}) \textrm{ for } i = 1, \ldots, n \]
\end{itemize}
\end{frame}

\begin{frame}
\frametitle{Independence Assumptions for Tagging}
\begin{block}{Chain Rule}
\[ P( t_{[1:n]} \mid x_{[1:n]}) = \pause \prod_{j=1}^n P(t_j \mid t_{j-1}, \ldots, t_1, x_{[1:n]}, j) \]
\end{block}
\pause
\begin{block}{Make independence assumptions}
\[ P( t_{[1:n]} \mid x_{[1:n]}) \approx \pause \prod_{j=1}^n P(t_j \mid t_{j-1}, x_{[1:n]}, j) \]
\centering {\small $j$ is the word being tagged. }\\
\centering {\color{red}\small We model the conditional probability directly: no Bayes Rule here. }
\end{block}
\pause
\begin{block}{Questions}
\begin{itemize}
\item Split up $P(t_j \mid t_{j-1}, x_{[1:n]}, j)$ into parameters?
\item How to find $\arg\max_{t_{[1:n]}} P( t_{[1:n]} \mid x_{[1:n]})$?
\end{itemize}
\end{block}
\end{frame}

\section{Log-linear models for Tagging}
\frame{\tableofcontents[currentsection]}

\begin{frame}
\frametitle{Representation: finding the right parameters}
\begin{block}{Problem: Predict ?? using context, $P(?? \mid \texttt{context})$ }
Profits\postag{N} soared\postag{V} at\postag{P} Boeing\postag{??} Co. , easily topping forecasts on Wall Street , as their CEO Alan Mulally announced first quarter results .
\end{block}
\pause
\begin{block}{Representation: history}
\begin{itemize}
\item A history is a 3-tuple: $(t_{-1}, x_{[1:n]}, i)$
\item $t_{-1}$ is the previous tag (we are assuming a bigram model)
\item $x_{[1:n]}$ are the $n$ words in the input
\item $i$ is the index of the word being tagged
\item For example, for $x_4$ = Boeing:
    \begin{itemize}
    \item $t_{-1}$ = P
    \item $x_{[1:n]}$ = (Profits, soared, ..., results, .)
    \item $i$ = 4
    \end{itemize}
\end{itemize}
\end{block}
\end{frame}

\newcommand{\featurefunction}[3]{
\[ 
f_{#1}(h,t) = \left\{
\begin{array}{ll}
1 & \textrm{if #2 and $t$ = #3 } \\
0 & \textrm{otherwise}
\end{array}
\right.
\]
}

\begin{frame}
\frametitle{Feature-vectors over history-tag pairs}
\begin{block}{Take a history, tag pair $(h,t)$}
$f_k(h, t)$ for $k = 1, \ldots, m$ are \textbf{feature functions} representing the tagging decision.
\end{block}
\pause
\begin{block}{Example: Part-of-speech tagging [Ratnaparkhi 1996]}
\centering
\featurefunction{100}{current word $x_i$ is \texttt{Boeing}}{\texttt{N}}
\featurefunction{101}{$t_{-1}$ is \texttt{P}}{\texttt{N}}
\end{block}
\end{frame}

\begin{frame}
\frametitle{Log linear model for Tagging}
\begin{itemize}[<+->]
\item Let there be $m$ features, $f_k(\textbf{x}, \textbf{y})$ for $k = 1, \ldots, m$
\item $\textbf{x} = x_{[1:n]}$ and $\textbf{y} = t_{[1:n]}$
\item Define a parameter vector $\textbf{w} \in \mathbb{R}^m$
\item Each $(\textbf{x}, \textbf{y})$ pair is mapped to score:
\[ s(\textbf{x}, \textbf{y}) = \sum_k w_k \cdot f_k(\textbf{x}, \textbf{y}) \]
\item Using inner product notation:
\begin{eqnarray*}
\textbf{w} \cdot \textbf{f}(\textbf{x}, \textbf{y}) & = & \sum_k w_k \cdot f_k(\textbf{x}, \textbf{y}) \\
s(\textbf{x}, \textbf{y}) & = & \textbf{w} \cdot \textbf{f}(\textbf{x}, \textbf{y}) 
\end{eqnarray*}
\item To get a probability from the score: Renormalize! 
\[ \Pr(\textbf{y} \mid \textbf{x}, \textbf{w}) = \frac{exp\left(s(\textbf{x}, \textbf{y})\right)}{\sum_{\textbf{y'}} exp\left(s(\textbf{x}, \textbf{y'})\right)} \]
\end{itemize}
\end{frame}

\begin{frame}
\frametitle{Feature functions for Tagging}
\begin{block}{Problem}
\begin{itemize}
\item We have defined a log-linear model using feature functions: $\textbf{f(x,y)}$ 
\item We have defined parameters using a history $h$ so feature functions are: $\textbf{f}(h, t)$
\end{itemize}
\end{block}
\end{frame}

\begin{frame}
\frametitle{Locally normalized log-linear taggers}
\begin{block}{Conditional Distribution over history, tag pair $(h,t)$}
\[ \log \Pr(t \mid h) = \textbf{w} \cdot \textbf{f}(h, t) - log \sum_{t'} \exp \left( \textbf{w} \cdot \textbf{f}(h, t') \right) \]
\begin{itemize}
\item $\textbf{f}(h, t)$ is a vector of feature functions
\item $\textbf{w}$ is the weight vector
\end{itemize}
\end{block}
\pause
\begin{block}{Local normalization for tagging}
\begin{itemize}
\item Word sequence: $x_{[1:n]}$ and tag sequence: $t_{[1:n]}$
\item Histories $h_i = (t_{i-1}, x_{[1:n]}, i)$ 
\end{itemize}
\[ \log \Pr( t_{[1:n]} \mid x_{[1:n]} ) = \sum_{i=1}^n \log \Pr(t_i \mid h_i) \]
\end{block}
\end{frame}


\begin{frame}
\frametitle{Globally normalized log-linear taggers}
\begin{block}{Global feature function $\mathbf{\Phi}(\textbf{x}, \textbf{y})$}
\begin{itemize}
\item Word sequence: $\textbf{x} = x_{[1:n]}$ and tag sequence: $\textbf{y} = t_{[1:n]}$
\item From \textit{local} histories $h_i = (t_{i-1}, x_{[1:n]}, i)$ to global $\Phi$ values:
\end{itemize}
\[ \Phi_k(x_{[1:n]}, t_{[1:n]}) = \sum_{i=1}^n f_k(h_i, t_i) \]
\begin{itemize}
\item $\mathbf{\Phi}(\textbf{x}, \textbf{y}) = (\Phi_1, \Phi_2, \ldots, \Phi_m)$ is a \textit{global} feature vector
\item $\textbf{w}$ is the weight vector for $\mathbf{\Phi}$
\end{itemize}
\end{block}
\pause
\begin{block}{Global normalization for tagging}
\[ \log \Pr(\textbf{y} \mid \textbf{x}, \textbf{w}) = \textbf{w} \cdot \mathbf{\Phi}(\textbf{x}, \textbf{y}) - \log \sum_{\textbf{y'}} exp\left( \textbf{w} \cdot \mathbf{\Phi}(\textbf{x}, \textbf{y'}) \right) \]
\end{block}
\end{frame}

\begin{frame}
\frametitle{Conditional Random Field}
\begin{block}{Global normalization for tagging}
\[ \log \Pr(\textbf{y} \mid \textbf{x}, \textbf{w}) = \textbf{w} \cdot \mathbf{\Phi}(\textbf{x}, \textbf{y}) - \log \sum_{\textbf{y'}} exp\left( \textbf{w} \cdot \mathbf{\Phi}(\textbf{x}, \textbf{y'}) \right) \]
\begin{itemize}
\item This model is also called a conditional random field (CRF)
\end{itemize}
\end{block}
\pause
\begin{block}{Algorithms for training and decoding}
\begin{itemize}
\item Global normalization could be expensive: requires sum over exponentially many terms $\textbf{y'}$
\item Finding $\arg\max_{\textbf{y}} \log \Pr(\textbf{y} \mid \textbf{x})$ can be accomplished using the Viterbi algorithm.
\item Training: finding the weight vector $\textbf{w}$ can be done using a variant of the Forward algorithm.
\end{itemize}
\end{block}
\end{frame}


\section*{Acknowledgements}

\begin{frame}
\centering
\begin{alertblock}{Acknowledgements}
Many slides borrowed or inspired from lecture notes by Michael Collins, Chris Dyer, Kevin Knight, Chris Manning, Philipp Koehn, Adam Lopez, Graham Neubig, Richard Socher and Luke Zettlemoyer from their NLP course materials. 

\bigskip

All mistakes are my own.

\bigskip

A big thank you to all the students who read through these notes and helped me improve them.

\end{alertblock}
\end{frame}



\end{document}

\begin{frame}
\frametitle{Maximum Entropy}
\begin{itemize}[<+->]
\item The log-linear model has an interpretation as a {\it maximum entropy} model.
\item For observed events, maximize likelihood. For unobserved events, from all
consistent models pick the one with maximum entropy.
\item The maximum entropy principle: related to Occam's razor and
  other similar justifications for scientific inquiry
\item Make the minimum possible assumptions about unseen data
\item Also: Laplace's {\em Principle of Insufficient Reason}: when one
  has no information to distinguish between the probability of two
  events, the best strategy is to consider them equally likely
\end{itemize}
\end{frame}

\begin{frame}
\frametitle{Logistic Regression}
\begin{itemize}
\item models effects of explanatory variables on binary valued
  variable
\item observations ${\bf x} = \{x_1, \ldots, x_j\}$ with success given
  by $q({\bf x})$: \[ q({\bf x}) = \frac{e^{g({\bf x})}}{1 + e^{g({\bf
  x})}} \] and \[ g({\bf x}) = \beta_0 + \sum_{j=1}^{k} \beta_j x_j \]
\end{itemize} 
\end{frame} 

\begin{frame}
\frametitle{Logistic Regression}
\begin{itemize}
\item probability that observations lead to success, or $p(a=1 \mid
  b)$:
 \[ p(a=1 \mid b) = \frac{e^{g(b)}}{1 + e^{g(b)}} \] where
 \[ g(b) = \beta_0 f_0(1,b) + \sum_{j=1}^{k} \beta_j f_j(1,b) \]
\item $\beta_j = \textrm{log} \alpha_j$, $f_0(1,b) = 1$ and $f_j(1,b) = x_j$
\end{itemize} 
\end{frame} 
